\chapter{Fallback Reconciler: Dynamic Structure When Blocks Aren't Enough}

Blocks are fast because structure is stable: the template doesn’t change, only slot values do.

But JSX in the real world has loops, conditionals, fragments, and dynamic lists.
HRBR needs a safe escape hatch for those regions.

That escape hatch is the \textbf{fallback reconciler}:

\begin{itemize}
  \item a runtime mount: \texttt{runtime/fallback.ts}
  \item a child-list reconciler: \texttt{runtime/reconciler.ts}
\end{itemize}

This chapter is grounded in:

\begin{itemize}
  \item \texttt{runtime/fallback.ts}
  \item \texttt{runtime/reconciler.ts}
  \item \texttt{runtime/__tests__/fallback.test.ts}
  \item \texttt{runtime/__tests__/reconciler.test.ts}
  \item \texttt{runtime/__tests__/reconciler.fuzz.test.ts}
\end{itemize}

\section{A tiny contract (what fallback + reconciler promise)}

The fallback reconciler is the HRBR answer to “dynamic structure”:

\begin{itemize}
  \item Blocks are optimal when structure is stable.
  \item Fallback reconciliation is the safe, general mechanism when structure changes.
\end{itemize}

The contract is intentionally narrow:

\begin{enumerate}
  \item \texttt{reconcileChildren(host, next)} only patches \textbf{direct children} of \texttt{host}.
  \item Each child spec is responsible for building/patching its own subtree (if any).
  \item In keyed mode, keys must be unique among siblings.
  \item In keyed mode, DOM identity is preserved by key when possible; the algorithm minimizes moves, but \emph{correctness} is the first goal.
\end{enumerate}

	exttt{mountFallback(host, render)} is the reactive wrapper:

\begin{itemize}
  \item \texttt{render()} can read signals.
  \item When those signals change, fallback schedules a reconcile pass in a scheduler lane.
  \item Disposal stops further reactive updates.
\end{itemize}

\section{Contracts}

\subsection{The reconciler contract: direct children only}

\texttt{reconcileChildren(host, next, options?)} patches \texttt{host.childNodes} so they match \texttt{next}.

Important: it reconciles \textbf{direct children only}. Each child spec is responsible for patching its own subtree.

That design choice keeps the reconciler small and composable:

\begin{itemize}
  \item A “child spec” can represent a component, a fragment, or even an embedded compiled block.
  \item Nested reconciliation happens inside that child spec’s own \texttt{patch()} function.
\end{itemize}

Types (\texttt{runtime/reconciler.ts}):

\begin{itemize}
  \item \texttt{ReconcileNode}: a logical node that maps to exactly one DOM node
  \item \texttt{ReconcileRange}: a logical item that maps to multiple DOM nodes
  \item \texttt{ReconcileItem = ReconcileNode | ReconcileRange}
\end{itemize}

Each item supports:

\begin{itemize}
  \item \texttt{create()} to create DOM
  \item \texttt{patch(existing)} to patch in place
  \item optional \texttt{destroy(existing)} to cleanup when removed
\end{itemize}

\subsection{Keyed vs unkeyed mode}

\texttt{reconcileChildren} chooses mode:

\begin{itemize}
  \item if \texttt{options.keyed} is provided, use it
  \item else infer keyed mode if any next item has a key
\end{itemize}

It also emits dev-only warnings for suspicious usage patterns when \texttt{opts.dev === true}:

\begin{itemize}
  \item mixing keyed and unkeyed siblings (unkeyed nodes become “unstable”)
  \item replacing an entire keyed list with a disjoint set of keys (often unstable keys like array index)
\end{itemize}

The two algorithms are:

\begin{itemize}
  \item \texttt{reconcileChildrenUnkeyed(host, next)}
  \item \texttt{reconcileChildrenKeyed(host, next, dev)}
\end{itemize}

\section{Why ranges exist (one logical item => multiple DOM nodes)}

In UI rendering, one logical item can expand to multiple nodes:

\begin{itemize}
  \item a fragment
  \item a component that returns multiple roots
  \item a compiled block or template that expands into siblings
\end{itemize}

HRBR models this with \texttt{ReconcileRange}:

\begin{verbatim}
{ kind: 'range', key?, create(): Node[], patch(nodes: Node[]) }
\end{verbatim}

The reconciler stores the range length on the first DOM node:

\begin{itemize}
  \item \texttt{__hrbrKey}: key of the logical item (on the start node)
  \item \texttt{__hrbrRangeLen}: number of DOM nodes in the range (on the start node)
\end{itemize}

This metadata is the core trick that lets the reconciler treat DOM as a list of \emph{logical items} rather than raw nodes.
The helper \texttt{buildCurrentRanges(host)} reads \texttt{childNodes} and groups them into ranges by repeatedly:

\begin{itemize}
  \item read \texttt{len = __hrbrRangeLen ?? 1} from the next start node
  \item slice \texttt{len} nodes and treat them as one logical item
  \item advance by \texttt{len}
\end{itemize}

This is also why correctness requires normalization when items change shape (node \(\leftrightarrow\) range).

Spec: \texttt{keyed: supports range items (logical item => multiple DOM nodes)}.

\section{Unkeyed reconciliation: simple and predictable}

\texttt{reconcileChildrenUnkeyed} relies on positional patching:

\begin{enumerate}
  \item patch common prefix (min(oldLen, newLen))
  \item remove extra old nodes
  \item append new nodes
\end{enumerate}

Spec: \texttt{unkeyed: patches in order and truncates/appends}.

Unkeyed mode is fast and simple, but it can’t preserve identity across reorders.

Unkeyed mode is best when:

\begin{itemize}
  \item lists are small
  \item items rarely reorder
  \item you don’t care about preserving DOM identity across reorder (e.g. no input state)
\end{itemize}

\section{Keyed reconciliation: preserve identity and minimize moves}

Keyed mode tries to:

\begin{itemize}
  \item reuse DOM nodes by key
  \item patch in place instead of recreating
  \item move existing ranges/nodes to new positions
  \item minimize DOM moves using a Longest Increasing Subsequence (LIS)
\end{itemize}

Specs in \texttt{runtime/__tests__/reconciler.test.ts} cover the basic expectations:

\begin{itemize}
  \item \texttt{keyed: moves DOM nodes instead of recreating them}
  \item \texttt{keyed: removes nodes that are no longer present}
  \item \texttt{keyed: inserts new nodes at the right position}
\end{itemize}

\subsection{Two-ended scan + LIS}

The reconciler uses a standard sequence:

\begin{enumerate}
  \item Build a map from key to current range (start..end, nodes[]).
  \item Fast two-ended scan for common prefix/suffix keys.
  \item For the middle region:
    \begin{itemize}
      \item build an "old index" sequence for new keys
      \item compute LIS indices
      \item move only the ranges not in the LIS
    \end{itemize}
\end{enumerate}

Why LIS?

\begin{itemize}
  \item If a subsequence of items already appears in increasing old-index order, those items can stay where they are.
  \item Everything else must move (or be created), but LIS keeps that set minimal.
  \item The implementation uses \texttt{lisIndices(seq)} over old indices, ignoring \texttt{-1} (new items).
\end{itemize}

Implementation cues:

\begin{itemize}
  \item \texttt{buildCurrentRanges(host)} groups DOM nodes into logical ranges by \texttt{__hrbrRangeLen}.
  \item \texttt{lisIndices(seq)} computes LIS over old indices (ignoring \texttt{-1}).
  \item \texttt{moveRangeBefore(host, range, anchor)} moves nodes end-to-start to preserve order.
\end{itemize}

\subsection{Dev diagnostics for suspicious usage}

Keyed mode includes two dev-only warnings:

\begin{itemize}
  \item Mixed keyed and unkeyed children in the same list.
  \item Next keys have no overlap with previous keys (often indicates unstable keys).
\end{itemize}

Specs:

\begin{itemize}
  \item \texttt{dev diagnostics: warns on mixed keyed/unkeyed children in keyed mode}
  \item \texttt{dev diagnostics: warns when next keys have no overlap with previous keys}
\end{itemize}

It also asserts that keys are unique among siblings.

That uniqueness assertion is a hard error, not a warning:

\begin{quote}
	exttt{assert(!seen.has(k), '[hrbr/reconciler] duplicate key ...')}
\end{quote}

This is intentional: duplicate keys make identity preservation ill-defined.

\subsection{Shape-change safety: node \(\leftrightarrow\) range}

A tricky edge case: a key might render as a singleton node in one render and as a range in another.

The reconciler handles this explicitly:

\begin{itemize}
  \item if key matches but logical shape changes (node vs range), it removes the old nodes and recreates
  \item it aggressively removes any \emph{unkeyed trailing nodes} that might remain after a range collapses into a singleton
\end{itemize}

This is where the metadata invariants matter. The code contains normalization passes that:

\begin{itemize}
  \item remove unkeyed nodes owned by keyed mode
  \item delete stale \texttt{__hrbrRangeLen} when the logical item is now a singleton
\end{itemize}

The fuzz test exists largely to guard these invariants.

One specific implementation detail to notice is \texttt{cleanupTrailingRangeNode(start, key)}.
When a keyed item collapses from a range into a singleton node, the DOM may still contain trailing unkeyed nodes that used to be part of that range.
The reconciler aggressively removes any contiguous unkeyed siblings immediately after the start node.

This is later reinforced by a final normalization pass that removes any unkeyed nodes interleaved between keyed logical items.

\section{Property/fuzz testing as an executable oracle}

\texttt{runtime/__tests__/reconciler.fuzz.test.ts} runs 250 seeded steps.

It compares \texttt{reconcileChildren(host, ...)} to a reference renderer that rebuilds the list each step but reuses nodes by key.

The key correctness check is not "minimal moves". It is:

\begin{quote}
For the same list of logical items, the resulting DOM represents the same logical structure and text content.
\end{quote}

The fuzz harness reads the logical structure with:

\begin{itemize}
  \item start key: \texttt{__hrbrKey}
  \item range length: \texttt{__hrbrRangeLen} (default 1)
\end{itemize}

Then it compares the logical tuples (key, len, texts[]) against the reference.

The fuzz test is doing something important for maintenance:

\begin{itemize}
  \item It doesn’t claim the reconciler is optimal.
  \item It claims that for any seeded sequence of updates, the reconciler produces the \emph{same logical output} as a simpler reference “rebuild” strategy.
\end{itemize}

That’s exactly the right kind of oracle for reconciliation code: optimality is a performance concern, but correctness is a binary property.

\section{Fallback mount: reconciling signal-driven dynamic children}

\texttt{runtime/fallback.ts} provides a reactive mount function:

\begin{verbatim}
mountFallback(host, render, { lane?, scheduler?, keyed? })
\end{verbatim}

Contract:

\begin{itemize}
  \item \texttt{render()} may read signals.
  \item When signals change, schedule a reconcile pass in the requested lane.
  \item The render result is \texttt{\{ children: ReconcileItem[], keyed? \}}.
\end{itemize}

Implementation uses two passes:

\begin{itemize}
  \item An effect that runs \texttt{render()} just to track dependencies, then schedules.
  \item A \texttt{runOnce()} pass that actually calls \texttt{reconcileChildren()}.
\end{itemize}

This split is subtle but intentional:

\begin{itemize}
  \item The tracking effect should be very cheap; it exists to subscribe to signals.
  \item The DOM work should be scheduled and batched (via \texttt{batch(() => runOnce())}).
  \item A \texttt{pending} flag ensures multiple invalidations coalesce into one scheduled reconcile.
\end{itemize}

In other words, \texttt{mountFallback} mirrors the same “coalesce and commit” philosophy used by \texttt{mountCompiledBlock} (Chapter 14): collect changes reactively, then perform one DOM commit.

Spec: \texttt{reconciles structural changes when signal-driven list changes} (\texttt{fallback.test.ts}).

That integration test verifies that:

\begin{itemize}
  \item reorders preserve node identity by key
  \item inserts create new nodes
  \item removals shrink the DOM
  \item disposing stops reactive updates
\end{itemize}

\section{Takeaways}

\begin{itemize}
  \item The fallback reconciler gives HRBR a safe strategy for dynamic structure.
  \item Keyed reconciliation preserves identity and minimizes moves via two-ended scan + LIS.
  \item Ranges let one logical item expand to multiple DOM nodes, enabling fragments and multi-root outputs.
  \item Metadata invariants (\texttt{__hrbrKey}, \texttt{__hrbrRangeLen}) are crucial; fuzz tests enforce them.
  \item \texttt{mountFallback} bridges signals/scheduler to \texttt{reconcileChildren} for real apps.
\end{itemize}

\section{Walking the tests (the executable spec)}

\subsection{\texttt{runtime/__tests__/reconciler.test.ts}}

This suite locks down the core expectations:

\begin{itemize}
  \item \textbf{Dev warnings are emitted:} by temporarily replacing \texttt{console.warn}, the tests prove warnings fire for
  mixed keyed/unkeyed children and for disjoint key sets.
  \item \textbf{Keyed reorder preserves identity:} the “moves DOM nodes” test captures original nodes \texttt{a0/b0/c0} then reorders and asserts node identity is preserved.
  \item \textbf{Keyed removals:} nodes that disappear from \texttt{next} are removed, while survivors keep identity.
  \item \textbf{Keyed inserts:} missing keys are created and inserted at the correct position.
  \item \textbf{Unkeyed mode is positional:} patches first, then truncates/appends, without identity guarantees across reorder.
  \item \textbf{Ranges work:} a range key produces two DOM nodes and behaves as a single logical item during moves/inserts/removals.
\end{itemize}

\subsection{\texttt{runtime/__tests__/reconciler.fuzz.test.ts}}

This is the long-term safety net.
It runs 250 seeded steps and compares the reconciler output against a reference renderer.

Key idea: the test reads \emph{logical structure} from the DOM via \texttt{__hrbrKey} and \texttt{__hrbrRangeLen}, then compares (key, len, per-node text) tuples.
If there’s a mismatch, it throws an error that includes seed, step, next specs, and DOM dumps to make replay/debugging practical.

\subsection{\texttt{runtime/__tests__/fallback.test.ts}}

This test shows the reconciler in its intended reactive setting:

\begin{itemize}
  \item a signal produces a keyed list
  \item changing the signal reorders, inserts, and removes
  \item identity is preserved for reused keys
  \item \texttt{dispose()} stops further reactive updates
\end{itemize}

The test uses \texttt{await new Promise((r) => setTimeout(r, 0))} to wait for the scheduled reconcile to run.

\section{Online resources}

\begin{itemize}
  \item A classic keyed diff discussion (and why identity matters):
  \href{https://react.dev/learn/rendering-lists}{React docs: Rendering Lists (keys)}.
  \item Longest Increasing Subsequence (the core move-minimization idea):
  \href{https://en.wikipedia.org/wiki/Longest_increasing_subsequence}{Wikipedia: Longest increasing subsequence}.
  \item A VDOM implementation reference with LIS-based keyed diff (for historical comparison):
  \href{https://vuejs.org/guide/essentials/list.html#maintaining-state-with-key}{Vue docs: Maintaining State with \texttt{key}}.
  \item DOM child manipulation primitives used by the reconciler:
  \href{https://developer.mozilla.org/en-US/docs/Web/API/Node/insertBefore}{MDN: \texttt{Node.insertBefore}} and
  \href{https://developer.mozilla.org/en-US/docs/Web/API/Node/removeChild}{MDN: \texttt{Node.removeChild}}.
\end{itemize}
